\section{Vascularization and Oxygen Dynamics in the Developing Cortex}

In addition to the generation of new cells in the neocortex, the developing brain must also establish a reliable supply of nutrients and oxygen. Vascularization begins in mice around E8.5–E10 with the formation of a vascular network surrounding the neural tube, called the perineural vascular plexus \citep{HOGANNeuralTubePatterns2004},
reviewed in \citep{JAMESNeuronalActionDeveloping2011}. From this outer network, new blood vessels grow radially from the pial toward the ventricular (apical) zone. Before reaching the inner (apical) surface of the neural tube, these vessels turn and extend sideways, where they connect with one another to form an internal vascular network, referred to as the subventricular vascular plexus \citep{TATAVascularisationCentralNervous2015, FANTINNRP1ActsCell2013}.
\\

Importantly, this process is not uniform across the CNS. While the hindbrain and spinal cord are vascularized relatively homogeneously, the forebrain exhibits a temporal delay, with blood vessels initially invading ventral regions while dorsal areas remain temporarily avascular \citep{VASUDEVANCompartmentspecificTranscriptionFactors2008}.
\\

Oxygen levels have been shown to play a critical role during neocortex development. Hypoxia during fetal development can disrupt cortical progenitor cell growth, leading to brain abnormalities such as smaller overall brain size, thinner cortex or enlarged ventricles \citep{piesovaImpactPerinatalHypoxia2020}.
In most studies, oxygen levels are quantified using  arterial partial oxygen tension ($PO_2$), where atmospheric oxygen levels correspond to a $PO_2$ of 21\% (150mmHg). 
\\

Even though hypoxia can disrupt normal brain development under pathological conditions, physiological hypoxia is essential in maintaining normal developmental processes. For instance, the stemness of NSCs is preserved under physiological hypoxic conditions \citep{EZASHILowO2Tensions2005}, aligning with observations that oxygen levels in the forebrain are significantly lower ($PO_2$ = 1-6\%) during development \citep{LEEDeterminationHypoxicRegion2001,ZHANGOxygenKeyFactor2011}.
Within this hypoxic environment, the SVZ reaches slightly higher $PO_2$ values (2.5-3\%) \citep{SANTILLIMildHypoxiaEnhances2010}.
\\

Similarly, the differentiation of NPCs into neurons is regulated by local oxygen levels \citep{malikNeurogenesisContinuesThird2013},  and increased oxygen availability has been associated with enhanced neuronal differentiation \citep{STORCHFunctionalCharacterizationDopaminergic2003,STUDEREnhancedProliferationSurvival2000,SANTILLIMildHypoxiaEnhances2010}.
\\

Consistent with these findings, NSC differentiation has been shown to coincide with blood vessel ingrowth \textit{in vivo} \citep{SHENEndothelialCellsStimulate2004},  
that is accompanied by increased tissue oxygenation and reduced levels of HIF-1$\alpha$, a key hypoxia marker. Additionally, evidence from genetic knockout supports a causal role of oxygen in cortical development. 
In Gpr124 knockout embryos, which exhibit impaired brain vascularization, oxygen levels fail to increase within the stem cell niche, resulting in sustained AP expansion at the expense of differentiation. Notably, experimentally increasing tissue oxygenation does not rescue vascular defects but is sufficient to reduce HIF-1$\alpha$ levels and restore progenitor differentiation \citep{LANGEReliefHypoxiaAngiogenesis2016}.
\\

While the temporal dynamics of vascularization during cortical development are relatively well understood, far less is known about the spatial distribution of oxygen within the developing cortex. Studies in the adult brain have revealed the presence of gradients, where oxygen levels decline along the length of capillaries due to consumption by surrounding cells. These gradients are particularly pronounced in grey matter, which is densely populated with neurons and exhibits a steeper oxygen drop-off compared to white matter, where myelinated axons reside. Notably, while grey matter shows sharper gradients, absolute $PO_2$
values tend to be lower in white matter \citep{erecinskaTissueOxygenTension2001}.
\\

In addition to the extracellular oxygen gradients, intracellular oxygen gradients are also assumed to exist in the brain, particularly in neurons, which maintain high rates of oxygen consumption. Within these cells, mitochondria may create localized regions of lower oxygen tension due to rapid oxygen utilization.
\\

Given that blood vessels invade the cortex radially from the pial surface toward deeper regions, it is reasonable to infer the existence of a spatial oxygen gradient. In such a scenario, more superficial regions populated by ULN and DLN would experience higher oxygen availability, whereas deeper regions including the VZ and SVZ that harbor neural stem and progenitor cells, would remain comparatively hypoxic. 
